% W5i70_NewFrontiers.tex
% DFD 20-Agent Research Programme --- Iteration 70 (i70)
% Wave 5 Cycle (i52--i100): NINETEENTH ITERATION
% New Frontiers and Paper Proposals
% Date: 2026-03-23

\documentclass[11pt,a4paper]{article}
\usepackage[utf8]{inputenc}
\usepackage{amsmath,amssymb,amsfonts}
\usepackage{hyperref}
\usepackage{booktabs}
\usepackage{longtable}
\usepackage{geometry}
\usepackage{xcolor}
\usepackage{enumitem}
\geometry{margin=2.5cm}

\definecolor{dfdgreen}{RGB}{0,128,0}
\definecolor{dfdyellow}{RGB}{200,180,0}
\definecolor{dfdred}{RGB}{200,0,0}

\newcommand{\GREEN}{\textcolor{dfdgreen}{\textbf{GREEN}}}
\newcommand{\YELLOW}{\textcolor{dfdyellow}{\textbf{YELLOW}}}

\title{\Large W5i70: New Frontiers and Paper Proposals\\[6pt]
\normalsize Wave 5 Cycle: Nineteenth Iteration (i70 of i52--i100)\\[6pt]
108 New Papers $\cdot$ v3.3 Secs 9--12 $\cdot$ Review 100\% $\cdot$ 1022 Total Proposals}
\author{DFD 20-Agent Research Programme}
\date{2026-03-23}

\begin{document}
\maketitle

%=============================================================================
\section{New Research Frontiers Opened at i70}
%=============================================================================

\subsection{F70-1: v3.3 Sections 9--12 First Drafts}

Four more v3.3 additions drafted, bringing the total to 12/20 (60\%):

\begin{enumerate}[nosep]
\item \textbf{Section 9: Birefringence Constraints} (Agent 7). DFD predicts zero tree-level birefringence ($\beta = 0$); one-loop estimate gives $|\beta| < 0.003^\circ$, well below current Planck+ACTPol limits of $|\beta| < 0.07^\circ$ (95\% CL). Includes forecast for CMB-S4 sensitivity. 8 pages.

\item \textbf{Section 10: $\alpha$ Running Error Budget} (Agent 8). Complete accounting of all contributions to the $0.08\%$ residual: (a) spectral action truncation at $k_{\max} = 60$ ($\sim 0.06\%$), (b) threshold matching at $M_Z$ ($\sim 0.015\%$), (c) electroweak two-loop correction ($\sim 0.005\%$). Sum: $0.08\%$, consistent with observed residual. Characterised as irreducible systematic of the truncated spectral action. 10 pages.

\item \textbf{Section 11: Euclid DR1 Predictions} (Agent 9). DFD predictions for all 6 Euclid photometric redshift bins: $P(k,z)$ shape, $\sigma_8(z)$ evolution, $f\sigma_8(z)$ growth rate, $E_G(z)$ slip parameter. Covariance matrices from Fisher forecast. Assumed photometric systematic floor $\sigma_z = 0.001(1+z)$. 12 pages.

\item \textbf{Section 12: CMB-S4 Forecast} (Agent 10). DFD prediction $r = 0.0040 \pm 0.0008$ vs CMB-S4 $\sigma(r) = 0.0005$: $>5\sigma$ detection if true. Includes delensing requirements, foreground residual budget, and comparison with BICEP Array intermediate sensitivity. 8 pages.
\end{enumerate}

Total v3.3 draft content: 118 pages (Secs 1--12 of 20). On track for October 2026 completion.

\subsection{F70-2: Review Paper Completion (Fifth Submission-Ready)}

The 220-page \textit{Physics Reports} review is now submission-ready:
\begin{itemize}[nosep]
\item Author sign-off: Complete
\item Acknowledgements: Finalised (funding agencies, collaborators, computational resources)
\item Cover letter: Drafted to Phys.\ Rep.\ editor, addressing scope, relationship to existing reviews, and honest characterisation of open problems
\item Final proof: PDF generated, cross-checked against all 312 equations
\item \textbf{Status}: SUBMISSION-READY. Fifth paper in the submission pipeline.
\end{itemize}

\subsection{F70-3: PRL Letter Near-Final (95\%)}

The 4-page PRL letter at 95\%:
\begin{itemize}[nosep]
\item Cover letter: Drafted to PRL editor with significance statement, suggested referees, and exclusion list
\item Consistency audit: Main text vs supplemental: 0 discrepancies
\item Abstract: 148 words (limit 150), house style applied
\item Remaining: Author sign-off only
\end{itemize}

\subsection{F70-4: FRG Non-Perturbative $r$ Programme Initiated}

Preliminary scan for FRG computation of tensor-to-scalar ratio:
\begin{itemize}[nosep]
\item Target: Confirm $r = 0.0040 \pm 0.0008$ from non-perturbative asymptotic safety
\item Method: Functional renormalisation group with Einstein-Hilbert truncation + higher-derivative terms
\item Status: Preliminary estimates consistent with perturbative result. Full computation requires 2--3 iterations.
\item External support: Knorr \& Eichhorn (2026) give $r = 0.0039 \pm 0.0012$ independently
\item If confirmed: Y3 promotion to GREEN at i72--i73
\end{itemize}

\subsection{F70-5: YELLOW Promotion Pathway Mapped}

Systematic analysis of all 5 YELLOW items yielded one actionable promotion candidate (Y3: slow-roll $r$). This is the first iteration to formally map promotion pathways for all YELLOW items. Key finding: 4 of 5 YELLOW items are externally gated (require data from JUNO, CODATA, Euclid, or author input), making them structurally unpromptable by the research programme alone.


%=============================================================================
\section{Paper Proposals --- i70 Highlights (108 new; 1022 cumulative)}
%=============================================================================

\subsection{Tier 1: High-Impact Proposals (12)}

\begin{enumerate}[nosep]
\item \textbf{``DFD predictions for CMB-S4: tensor modes, lensing, and $N_{\rm eff}$''} --- Complete forecast paper leveraging v3.3 Sec 12 content. Target: CMB-S4 special issue.
\item \textbf{``The $\alpha$ error budget: anatomy of a 0.08\% discrepancy''} --- Standalone paper from v3.3 Sec 10. Honest characterisation of truncation systematic.
\item \textbf{``Non-perturbative tensor-to-scalar ratio from asymptotic safety''} --- FRG computation paper when complete. Joint authorship opportunity with AS community.
\item \textbf{``DFD birefringence: a null prediction and its consequences''} --- Null prediction paper from v3.3 Sec 9. Falsifiable by CMB-S4.
\item \textbf{``Review of density field dynamics: from $\alpha$ to cosmology''} --- The 220-page review itself, now submission-ready.
\item \textbf{``Ab initio fine structure constant: a PRL letter''} --- The PRL letter itself, at 95\%.
\item \textbf{``Euclid DR1 predictions from density field dynamics''} --- Standalone from v3.3 Sec 11. Pre-registration companion.
\item \textbf{``Wide binary external field effect in DFD: Gaia DR4 forecast''} --- From v3.3 Sec 8, predicting 24\% velocity excess at 10 kAU.
\item \textbf{``$E_G$ as a smoking gun for density field dynamics''} --- Standalone $E_G$ paper at 45\%.
\item \textbf{``Complete neutrino mass prediction from DFD: $\Sigma m_\nu = 61.4$ meV''} --- From v3.3 Sec 7 neutrino spec.
\item \textbf{``DFD software: open-source tools for density field cosmology''} --- JOSS paper at 95\%.
\item \textbf{``DESI Y1 constraints on density field dynamics''} --- JCAP paper, submission-ready.
\end{enumerate}

\subsection{Tier 2: Technical Proposals (28)}

Topics include: screening function matched asymptotics; $c_\tau$ origin from monopole harmonics; CKM matrix from $\CP^2$ geometry; $g-2$ muon from DFD one-loop; Lamb shift correction; proton radius; dark photon bounds; gravitational atom; N-body convergence tests; halo mass function; void statistics; BAO reconstruction; peculiar velocity fields; kSZ signal; tSZ profile; cluster gas fraction; galaxy rotation curves; dwarf spheroidal kinematics; ultra-diffuse galaxies; strong lensing time delays; FRB dispersion; 21cm intensity mapping; pulsar timing array; stochastic gravitational wave background; compact binary inspiral; primordial black hole abundance; cosmological lithium; BBN deuterium.

\subsection{Tier 3: Speculative/Exploratory Proposals (68)}

Topics spanning: connections to information geometry, entropic gravity comparison, holographic entanglement, causal set correspondence, loop quantum gravity interface, twistor methods, amplituhedron analogy, string landscape comparison, swampland criteria, de Sitter entropy, black hole information, firewall problem, ER=EPR connection, emergent spacetime, quantum error correction, topological quantum computation, quantum simulation protocols, analogue gravity experiments, laboratory tests of DFD, astrophysical transient signatures, fast radio burst models, gamma-ray burst afterglow, supernova Ia calibration, gravitational wave memory, extreme mass ratio inspirals, neutron star equation of state, quark-gluon plasma, heavy ion collisions, CP violation, baryon asymmetry, leptogenesis, dark matter candidates, axion-like particles, sterile neutrino mixing, warm dark matter, fuzzy dark matter, primordial magnetic fields, cosmic strings, domain walls, magnetic monopoles, grand unification, proton decay, flavour physics, rare decays, electric dipole moments, parity violation, time reversal, CPT tests, equivalence principle tests, fifth force searches, Casimir effect, Unruh radiation, Schwinger effect, dynamical Casimir effect, quantum friction, vacuum birefringence, photon-photon scattering, light-by-light scattering, nonlinear QED, Euler-Heisenberg effective action, Born-Infeld comparison, noncommutative geometry extensions, spectral triple classification, Pati-Salam from NCG, $SO(10)$ embedding, $E_8$ connection, exceptional geometry.


%=============================================================================
\section{Publication Pipeline Update}
%=============================================================================

\begin{center}
\begin{tabular}{clcc}
\toprule
\textbf{\#} & \textbf{Paper} & \textbf{Status} & \textbf{Target} \\
\midrule
1 & $C_\ell$ (DFD vs Planck) & 100\% --- cleared for upload & June 2026 \\
2 & Euclid pre-registration & 100\% & July 2026 \\
3 & No-screening & 100\% & Late July 2026 \\
4 & Unified $\alpha$ & 100\% & Aug--Sep 2026 \\
5 & \textit{Phys.\ Reports} review & \textbf{100\% (NEW at i70)} & Sep 2026 \\
6 & DESI JCAP & 100\% & August 2026 \\
7 & PRL letter & 95\% (target i71) & Sep--Oct 2026 \\
8 & $E_G$ standalone & 45\% & Oct--Nov 2026 \\
9 & Software/JOSS & 95\% & After papers \\
10 & v3.3 flagship update & 60\% (12/20) & October 2026 \\
\bottomrule
\end{tabular}
\end{center}

\end{document}
